% Generated from main.typ. Typst is authoritative; regeneration overwrites this file.
\documentclass{qlnotes-export}

\title{MATH 597: Measure Theory}
\subtitle{Typst-first course notes and worked homeworks}
\author{Qiulin Fan}
\date{Winter 2025}

\begin{document}
\mainmatter
\chapter{\texorpdfstring{\(L_{p}\) space and inequalities}{L\_\{p\} space and inequalities}}\label{l_p-space-and-inequalities}

\section{\texorpdfstring{Banach Space and \(L_{p}\) space {[}Fol 5.1; 6.1{]}}{Banach Space and L\_\{p\} space {[}Fol 5.1; 6.1{]}}}

对应 Folland 5.1(1), 6.1(1).

\subsection{norm and completeness}\label{norm-and-completeness}

Recall:

\begin{definition}{\protect\phantomsection\label{kn-semi-norm-norm}{\textbf{semi-norm,
norm}}}

一个\textbf{semi norm} 是一个函数 \(\left. | \middle| \cdot \middle| \middle| :V\rightarrow\lbrack 0,\infty) \right.\) starting from a vector space \(V\). 其满足 (1): tri eq 和 (2): homogeneity.\\
如果一个 semi-norm 满足 (3): \(\left. | \middle| v \middle| \middle| = 0 \right.\) iff \(v = 0\), 则称它为一个 \textbf{norm}.

\end{definition}

\begin{definition}{\protect\phantomsection\label{kn-banach-space}{\textbf{Banach space}}}

一个 normed vector space \(\left. (V, \middle| \middle| \cdot \middle| \middle| ) \right.\) 的 induced metric space 如果是 complete 的, 它就被称为一个 \textbf{Banach space}.\\

\end{definition}

\begin{qlremark}{Remark}

Cauchy 指的是对于任意 \(\epsilon\), 都存在 \(N\) 使得对于任意 \(n,m \geq N\) 都有

\[\parallel v_{n} - v_{m} \parallel < \epsilon\]

而 convergent 指的是存在一个极限 \(v\), 使得对于任意 \(\epsilon\), 都存在 \(N\) 使得对于任意 \(n \geq N\) 都有

\[\parallel v_{n} - v \parallel \leq \epsilon\]

By tri ineq 容易证明: 在 genral normed VS 中, convergent imply Cauchy, 反之未必. convergent 是更强的条件. (interestingly, convergence in measure 却不 imply Cauchy in measure)

\end{qlremark}

\begin{example}

\({\mathbb{R}}^{n},{\mathbb{C}}^{n}\) with Euclidean norm is a Banach space.\\
\(C^{0}(\lbrack 0,1\rbrack)\): space of ctn functions on \(\lbrack 0,1\rbrack\) equipped with \(\sup\) norm \textbf{is Banach}.

\[\left. | \middle| f - g \middle| \middle| := \sup\limits_{x \in \lbrack 0,1\rbrack} \middle| f(x) - g(x)| \right.\]

\(C_{c}^{0}({\mathbb{R}})\): space of ctn functions with cpt supp on \(\mathbb{R}\) equipped with \(\sup\) norm \textbf{is not Banach}! 这是因为, 一个有 cpt supp 的 function seq 的极限未必有 cpt supp. 比如 \((\chi_{\lbrack - n,n\rbrack})_{n \in {\mathbb{N}}}\).

\end{example}

\begin{lemma}

A metric space \((X,\rho)\) is \textbf{complete} iff \textbf{every Cauchy seq has a subseq that converges.}

\end{lemma}

\begin{proof}

Trivial.\\
\(\Longrightarrow\): Clear.\\
\(\Longleftarrow\): subseq conv dist bound + Cauchy dist bound can bound the whole tail with arbitrary \(\epsilon\).

\end{proof}

这个 statement, 直接把 complete 的定义从每个 Cauchy seq 都收敛, 优化为每个 Cauchy seq 都有一个收敛 subseq.

\subsection{\texorpdfstring{every Cachy seq conv (complete) \(\Leftrightarrow\) every abs conv series convs}{every Cachy seq conv (complete) \textbackslash Leftrightarrow every abs conv series convs}}\label{every-cachy-seq-conv-complete-leftrightarrow-every-abs-conv-series-convs}

\begin{definition}{\protect\phantomsection\label{kn-series-convergence-absolute-convergence}{\textbf{series:
convergence 和 absolute convergence}}}

对于一个 normed VS \(\left. (V, \middle| \middle| \cdot \middle| \middle| ) \right.\) 中的 seq \((v_{n})\), 我们称 \(\sum_{n = 1}^{\infty}v_{n}\) \textbf{converges}, 如果存在 \(v \in V\) s.t.

\[\lim\limits_{N\rightarrow\infty}\sum\limits_{n = 1}^{N}v_{n} = v\]

即

\[\lim\limits_{N\rightarrow\infty} \parallel v - \sum\limits_{n = 1}^{N}v_{n} \parallel = 0\]

我们称 \(\sum_{n = 1}^{\infty}v_{n}\) \textbf{absolutely converges}, 如果

\[\left. \sum\limits_{n = 1}^{\infty} \middle| \middle| v_{n} \middle| \middle| < \infty \right.\]

即这个 series 对应的 norm series converges to some real number.

\end{definition}

% qlnotes: kind=theorem; id=thm-09-l-p-space-and-inequalities-another-criterion-for-banach-space
\begin{theorem}{\protect\phantomsection\label{kn-another-criterion-for-banach-space}{\textbf{another
criterion for Banach space}}}\label{thm-09-l-p-space-and-inequalities-another-criterion-for-banach-space}

A normed VS \(\left. (V, \middle| \middle| \cdot \middle| \middle| ) \right.\) is a Banach space iff every absolutely convergent series converges.

\end{theorem}

\begin{proof}

``\(\Longrightarrow\)": 如果 \(\left. (V, \middle| \middle| \cdot \middle| \middle| ) \right.\) is a Banach space, Suppose \(\left. \sum_{n = 1}^{\infty} \middle| \middle| v_{n} \middle| \middle| < \infty \right.\), 取部分和序列

\[S_{N} := \sum\limits_{n = 1}^{N}v_{n}\]

有

\[\parallel S_{m} - S_{n} \parallel = \parallel \sum\limits_{k = n + 1}^{m}v_{k} \parallel \leq \sum\limits_{k = n + 1}^{m} \parallel v_{k} \parallel\]

For large enough \(m,n\) 这个 bound 可以无限小, 因而 \((S_{N})\) is Cauchy. ``\(\Longleftarrow\)": 如果 \(\left. (V, \middle| \middle| \cdot \middle| \middle| ) \right.\) 中 every absolutely convergent series converges.\\
Suppose \((v_{n})\) is Cauchy. WTS it converges.\\
By Cauchy, 存在 subseq, say labeled \(n_{1} < n_{2} < \cdots\), s.t. \(\left. | \middle| v_{m} - v_{n} \middle| \middle| < \frac{1}{3^{j}} \right.\)for all \(m,n \geq n_{j}\) Then

\[\left. \sum\limits_{j = 1}^{\infty} \middle| \middle| v_{n_{j + 1}} - v_{n_{j}} \middle| \middle| < \infty \right.\]

Let \((y_{j})\) be s.t. \(y_{1} = v_{n_{1}}\), \(y_{j} = v_{n_{j + 1}} - v_{n_{j}}\), then

\[\sum\limits_{j = 1}^{\infty} \parallel y_{j} \parallel \leq \parallel y_{1} \parallel + \sum\limits_{j}\frac{1}{2^{j}} = \parallel y_{1} \parallel + 1 < \infty\]

并且有:

\[v_{n_{j}} = \sum\limits_{k = 1}^{j}y_{k}\]

由于 \(\sum_{j = 1}^{\infty} \parallel y_{j} \parallel < \infty\), by our assumption 得到, 这个极限 \(\lim_{j\rightarrow\infty}v_{n_{j}} = \sum_{k = 1}^{\infty}y_{k}\) 是存在的.

\end{proof}

\begin{qlremark}{Remark}

这个证明中也有一个简略但是有用的结论: 任意 normed VS 中, \textbf{一个 series absolutely convergent 可以推出它的部分和 seq 是 Cauchy 的. (反向则未必成立).}\\
整个 imply 关系的示意图:

\[\begin{matrix}
{\sum\limits_{k = 1}^{\infty} \parallel x_{k} \parallel < \infty\Longrightarrow S_{N}\ \text{Cauchy}} & {\overset{\text{if Banach}}{\Longrightarrow}S_{N}\ \text{converges}\Leftrightarrow\sum\limits_{k = 1}^{\infty}x_{k}\ \text{converges}\Longrightarrow(x_{k})\rightarrow 0} \\
 & \overset{\text{always}}{\Longleftarrow}
\end{matrix}\]

(这个图直观说明了为什么 Banach 和 "every abs conv seq conv" 是等价的. 因为这只是\textbf{在 Cauchy imply conv 的前后套了两个必然发生的 implication 关系}而已. 但有时候, 这个关系反而更加好证明.)\\
\textbf{(注意, partial sum seq Cauchy 并不 imply 原 series absolutely converge!})

\end{qlremark}

\subsection{任何 finite dim normed VS 一定 Banach, infinite dim 则不一定 Banach}\label{ux4efbux4f55-finite-dim-normed-vs-ux4e00ux5b9a-banach-infinite-dim-ux5219ux4e0dux4e00ux5b9a-banach}

\begin{qlremark}{Remark}

Note, 我们知道在 \({\mathbb{R}}^{n}\), \({\mathbb{C}}^{n}\) 上, abs conv 一定 imply con; 但是在 general (infinite dimension) 的 normed VS 上, \textbf{absolutely converge 并不 imply converge.}\\
1. As is known to all, \({\mathbb{R}}^{n},{\mathbb{C}}^{n}\) 上 Euclidean norm 的 induced metric 就是 Euclidean metric, making it complete metric space, 从而是 Banach space.\\
2. recall in elementary functional analysis:

\begin{definition}

我们称两个 norms \(\parallel \cdot \underset{a}{\parallel}, \parallel \cdot \underset{b}{\parallel}\) on a vector space 是 equivalent, 如果存在常数 \(C_{1},C_{2} > 0\) 使得对于任意 \(x\) 都有

\[C_{1} \parallel x\underset{a}{\parallel} \leq \parallel x\underset{b}{\parallel} \leq C_{2} \parallel x\underset{a}{\parallel}\]

这一定义即 topologically equivalent. 因为 equivalent norms define \textbf{equivalent metric, 从而 same topology.}

\end{definition}

以及这个经典的定理:

\begin{theorem}

finite dimensional vector space \(X\) 上, 所有 norms 都 equivalent.

\end{theorem}

这里先不证明.\\
利用这个定理, 我们发现 \textbf{\({\mathbb{R}}^{n},{\mathbb{C}}^{n}\) 上采用任何 norm 都是 Banach space.}\\
3. 我们 recall: 任何 finite dim \(\mathbb{R}\)-vector space 或者 \(\mathbb{C}\)-vector space 都 isomorphic to some \({\mathbb{R}}^{n}\), \({\mathbb{C}}^{n}\). 因而 利用这 theorem, 我们得到: \textbf{任何 finite dim normed VS 都是 complete metric space (Banach space), regardless of choice of norm.}\\
然而 \textbf{infinite dim normed VS 则未必一定 Banach.}\\
一个常见的反例:

\[V = {\mathbb{R}}\lbrack x\rbrack\]

所有的 polynomials with real coeffs. 考虑这一 norm:

\[\left. \parallel p\underset{\infty}{\parallel} = \sup\limits_{x \in \lbrack 0,1\rbrack} \middle| p(x)| \right.\]

\({\mathbb{R}}\lbrack x\rbrack\) 是无限维的, 因为多项式的次数可以任意提高.\\
\(({\mathbb{R}}\lbrack x\rbrack, \parallel p\underset{\sup}{\parallel})\) 不是 complete 的, 其 completion 是 Banach 空间 \(C\lbrack 0,1\rbrack\), 所有在 \(\lbrack 0,1\rbrack\) 上的连续函数, with the same \(\sup\) norm.\\

\end{qlremark}

下面我们将介绍一类 infinite dimension 但是 Banach 的 normed VS: \(L^{p}\) spaces.

\subsection{\texorpdfstring{\(L^{p}\) spaces}{L\^{}\{p\} spaces}}\label{lp-spaces}

\begin{definition}{\protect\phantomsection\label{kn-l-p-spaces}{\textbf{\(L_{p}\) spaces}}}

Consider \(p \in (0,\infty)\).\\
Let \((X,\mathcal{A},\mu)\) 为一个 measure space.\\
Define for \(f:X\rightarrow{\mathbb{R}}\) measurable:

\[\left. | \middle| f \middle| \middle| {}_{p}: = (\int \middle| f \middle| {}_{p} d\mu)^{\frac{1}{p}} \in \lbrack 0,\infty\rbrack \right.\]

Define

\[L^{p}(\mu): = \left\{ f: \middle| \middle| f \middle| \middle| {}_{p} < \infty \right\}/ \sim\]

where \(f \sim g\) if \(f = g\) a.e.

\end{definition}

固定一个 measure space \((X,\mathcal{A},\mu)\), 我们将用 \(L_{p}\) 来简易指代 \(L^{p}(\mu)\).\\

\begin{qlremark}{Remark}

注意, 我们容易发现:if \(0 < p < \infty\) and \(f\) measurable, TFAE：

\begin{itemize}
\item
  \(f \in L^{p}\)
\item
  \(\left. |f \middle| \in L^{p} \right.\)
\item
  \(\left. |f \middle| {}_{p} \in L^{1} \right.\)
\end{itemize}

\end{qlremark}

\begin{example}

\((X,\mathcal{A},\mu) := ({\mathbb{R}},\mathcal{L},m)\),

\[f(x): = \frac{1}{x^{\alpha}}\chi_{(0,1)}, f \in L^{p}(m)\Leftrightarrow\alpha p < 1\]\[f(x): = \frac{1}{x^{\alpha}}\chi_{(1,\infty)}, f \in L^{p}(m)\Leftrightarrow\alpha p > 1\]

\((X,\mathcal{A},\mu) := ({\mathbb{N}},\mathcal{P}({\mathbb{N}}),\mu_{counting})\),

\[L^{p}(\mu_{counting}) = \left\{ (a_{n})_{n \in {\mathbb{N}}}:\sum\limits_{n = 1}^{\infty} \middle| a_{n} \middle| {}_{p} < \infty \right\}\]

\end{example}

\begin{lemma}{\protect\phantomsection\label{kn-l-p-space-is-a-vector-space}{\textbf{\(L_{p}\)
space is a vector space}}}

\(L_{p}\) space is a \(\mathbb{C}\)-vector space.

\end{lemma}

\begin{proof}

Suppose \(f,g \in L^{p}\).\\
由于

\[\left. |f + g \middle| {}_{p} \leq ( \middle| f \middle| + \middle| g \middle| )^{p} \leq (2\max\left\{ |f \middle| , \middle| g| \right\})^{p} \leq 2^{p}( \middle| f \middle| {}_{p} + \middle| g \middle| {}_{p}) \right.\]

于是 by linearity of integral, 得到:

\[f,g \in L^{p}\Longrightarrow f + g \in L^{p}\]

(Note: \(p > 1\) 时也可以 by \(|x|^{p}\) 这一函数的 convexity 得到这个 bound, 但是这个方法只有效于 \(p > 1\))

\end{proof}

但是 \textbf{Question 1:} \textbf{Is \(L_{p}\) a normed VS? 即, \(\parallel \cdot \underset{p}{\parallel}\) 总是一个 valid norm 吗?} A: \textbf{True for \(p \in \lbrack 1,\infty)\), false for \(p \in (0,1)\).} Homogeneity 和 \(\parallel f\underset{p}{\parallel} = 0\) iff \(f = 0\) (a.e.) 是显然的, 但是我们发现, tri ineq 没有显然的证明.\\
Next lecture, we will show the Minkowski's ineq, 即 \(L^{p}\) space 上的三角不等式:

\[\left. | \middle| f + g \middle| \middle| {}_{p} \leq \middle| \middle| f \middle| \middle| {}_{p} + \middle| \middle| g \middle| |_{p} \right.\]

但是这个不等式只 hold for \(p \in \lbrack 1,\infty)\), 并且 fail otherwise.\\
(因而对于 \(L^{p}\) space 的研究, 我们将 \textbf{focus on \(p \in \lbrack 1,\infty)\) 的情况.})\\
\textbf{Question 2: Is \(L^{p}\) space, \(p \in \lbrack 1,\infty)\), Banach? Answer: Yes.}\\
我们也将在 next lecture 证明它.

\section{\texorpdfstring{inequilities on \(L^{p}\) spaces {[}Fol 6.1{]}}{inequilities on L\^{}\{p\} spaces {[}Fol 6.1{]}}}

对应 Folland 6.1(2).\\
我们将证明 Hölder's ineq 以及它的 corollary Minkowski's ineq, 从而证明: \(L^{p}\) 是一个 normed VS, 并且是一个 Banach space (这里 \(1 \leq p < \infty\), 但是 later we will also prove \(L^{\infty}\) 也是 Banach space).\\
这两个不等式非常重要.

\subsection{Hölder's ineq}\label{huxf6lders-ineq}

\begin{theorem}{\protect\phantomsection\label{kn-holder-inequality}{\textbf{Hölder's
ineq}}}

Consider conjugate pair: \(p,q \in \lbrack 1,\infty)\) s.t.

\[\frac{1}{p} + \frac{1}{q} = 1\]

则对于任意两个 measurable function \(f,g:X\rightarrow{\mathbb{C}}\), 一定有:

\[\parallel fg\underset{1}{\parallel} \leq \parallel f\underset{p}{\parallel} \cdot \parallel g\underset{q}{\parallel}\]

特别地, 如果 \(f \in L^{p}(\mu)\), \(g \in L^{q}(\mu)\), 则 \(fg \in L^{1}(\mu)\), 并且 equality holds iff

\[\left. \parallel g\underset{q}{\overset{q}{\parallel}} \middle| f \middle| {}_{p} = \parallel f\underset{p}{\overset{p}{\parallel}} \middle| g \middle| {}_{q}\quad\mu\ \text{-a.e.} \right.\]

\end{theorem}

\begin{qlremark}{Remark}

For \((p,q) = (2,2)\), this is \textbf{Cauchy-Swartz ineq}:

\[\parallel fg\underset{1}{\parallel} \leq \parallel f\underset{2}{\parallel} \cdot \parallel g\underset{2}{\parallel}\]

即:

\[\left. (\int f\,\bar{g}) \leq \int \middle| fg \middle| \leq \sqrt{\left. (\int \middle| f \middle| {}_{2})(\int \middle| g \middle| {}_{2}) \right.} \right.\]

\end{qlremark}

\begin{proof}

Trivial Case 1: 如果 \(\parallel f\underset{p}{\parallel} = 0\) (或者\(\parallel g\underset{q}{\parallel} = 0\) ), then \(f\) is zero \(\mu\)-almost everywhere, and the product \(fg\) is zero \(\mu\)-almost everywhere, 于是两边都是 \(0\), ineq trivially true.\\
Trivial Case 2: 如果 \(\parallel f\underset{p}{\parallel} = \infty\) or \(\parallel g\underset{q}{\parallel} = \infty\), 则右边 infinite, ineq trivially true. 因而我们只需要考虑 \(\parallel f\underset{p}{\parallel}\) and \(\parallel g\underset{q}{\parallel}\) are in \((0,\infty)\) 的情况就好了.\\
Main case: 我们需要一个 Lemma:

\begin{lemma}{\protect\phantomsection\label{kn-young-s-inequality-for-products}{\textbf{Young's
inequality for products}}}

Whenever \(p,q \in (1,\infty)\) with \(\frac{1}{p} + \frac{1}{q} = 1\), 都有

\[ab \leq \frac{a^{p}}{p} + \frac{b^{q}}{q},\quad\forall a,b \geq 0\]

where equality is achieved if and only if \(a^{p} = b^{q}\).\\
另一个等价形式是:

\[a^{\lambda}b^{1 - \lambda} \leq \lambda a + (1 - \lambda)b,\quad\forall a,b \geq 0\]

\end{lemma}

\begin{proof}

\textbf{of Lemma:}\\
\(b = 0\) 则 trivial case. 因而 setting \(t: = \frac{a}{b}\), reduced to show:

\[t^{\lambda} \leq \lambda t + (1 - \lambda)\]

with eq iff \(t = 1\). 这是显然的, 因为 by Calculus, \(t^{\lambda} - \lambda t\) 是 strictly increasing for \(t < 1\), strictly decreasing for \(t > 1\) 的, max 在 \(t = 1\), 正好是 \(1 - \lambda\).

\end{proof}

使用 Young's inequality for products 得到:

\[\frac{|f(x)|}{\parallel f\underset{p}{\parallel}}\frac{|g(x)|}{\parallel g\underset{q}{\parallel}} \leq \frac{|f(x)|^{p}}{p \parallel f\underset{p}{\overset{p}{\parallel}}} + \frac{|g(x)|^{q}}{q \parallel g\underset{q}{\overset{q}{\parallel}}},\quad x \in X\]

Integrating both sides gives

\[\frac{\parallel fg\underset{1}{\parallel}}{\parallel f\underset{p}{\parallel} \parallel g\underset{q}{\parallel}} \leq \frac{\parallel f\underset{p}{\overset{p}{\parallel}}}{p \parallel f\underset{p}{\overset{p}{\parallel}}} + \frac{\parallel g\underset{q}{\overset{q}{\parallel}}}{q \parallel g\underset{q}{\overset{q}{\parallel}}} = \frac{1}{p} + \frac{1}{q} = 1,\]

which proves the claim.\\
Integration 的 equality holds iff point equality holds a.e., 并且, by Young's inequality for products, 上面的 equality holds iff

\[\left. \parallel g\underset{q}{\overset{q}{\parallel}} \middle| f \middle| {}_{p} = \parallel f\underset{p}{\overset{p}{\parallel}} \middle| g \middle| {}_{q}\quad\mu\ \text{-a.e.} \right.\]

\end{proof}

\begin{qlremark}{Remark}

1. 显然, 根据我们的证明过程可知: \textbf{Hölder's ineq also holds on any measurable subset \(S \subset X\)}:

\[\left. \int_{S} \middle| fg \middle| \leq (\int_{S} \middle| f \middle| {}_{p})^{\frac{1}{p}}(\int_{S} \middle| g \middle| {}_{q})^{\frac{1}{q}} \right.\]

2. 这里的满足 \(\frac{1}{p} + \frac{1}{q} = 1\) 的 \(p,q\) 我们称之为: \textbf{Hölder conjugate}, 并称它们互为对方的 \textbf{conjugate exponent}.\\
3. 左边实际上是两个正值函数的 inner product, 相当于把一个投影到另一个上;\\
几何直观: Hölder's ineq 在退化为 Cauchy-Swartz 时表示, 两个函数/向量的内积一定小于等于长度积; 而 Hölder's ineq 更广义: 表示它们的内积一定小于它们取任意相互 conjugate 的 norm 长度的积. 并且 sooner 我们会学到: 对作为 Hölder conjugates 的 \(p,q\), \(L^{p}\) 和 \(L^{q}\) 互为 dual space, 从而 Hölder ineq 表示的是就是 norm 与其 dual norm 之间的 maximal inner product 控制关系.

\end{qlremark}

\begin{qlremark}{Remark}

Hölder's ineq 有一个 generalization: 对于任意 \(0 < s < \infty\) and \(0 < p_{1},\ldots,p_{n} < \infty\) such that

\[\frac{1}{p_{1}} + \frac{1}{p_{2}} + \ldots + \frac{1}{p_{n}} = \frac{1}{s};\]

都有

\[\parallel f_{1}f_{2}\cdots f_{n}\underset{s}{\parallel} \leq \parallel f_{1}\underset{p_{1}}{\parallel} \parallel f_{2}\underset{p_{2}}{\parallel}\cdots \parallel f_{n}\underset{p_{n}}{\parallel}.\]

This generalization will be proved in hw8.

\end{qlremark}

\subsection{\texorpdfstring{Minkowski's ineq: tri ineq on \(L^{p}\), 确认 \(\parallel \cdot \underset{p}{\parallel}\)-norm 是 \(L^{p}\) 上的 valid norm}{Minkowski's ineq: tri ineq on L\^{}\{p\}, 确认 \textbackslash parallel \textbackslash cdot \textbackslash underset\{p\}\{\textbackslash parallel\}-norm 是 L\^{}\{p\} 上的 valid norm}}\label{minkowskis-ineq-tri-ineq-on-lp-ux786eux8ba4-parallel-cdot-undersetpparallel-norm-ux662f-lp-ux4e0aux7684-valid-norm}

Minkowski's ineq 即 \(L^{p}\) space 上的 tri ineq.

\begin{corollary}{\protect\phantomsection\label{kn-minkowski-inequality}{\textbf{Minkowski
inequality}}}

对于任意 \(1 \leq p < \infty\), 都有:

\[\parallel f + g \parallel \leq \parallel f\underset{p}{\parallel} + \parallel g\underset{p}{\parallel}\]

\end{corollary}

\begin{proof}

显然, 对于任意 \(x\) 都有:

\[\left. |f + g \middle| {}_{p} \leq ( \middle| f \middle| + \middle| g \middle| ) \middle| f + g|^{p - 1} \right.\]

因而:

\[\left. \int \middle| f + g \middle| {}_{p} \leq \int \middle| f \middle| \cdot \middle| f + g \middle| {}_{p - 1} + \int \middle| g \middle| \cdot \middle| f + g|^{p - 1} \right.\]

我们定义

\[\left. h(x) := \middle| f(x) + g(x)|^{p - 1} \right.\]

于是

\[\begin{matrix}
\left. \int \middle| f + g|^{p} \right. & \left. \leq \int \middle| fh \middle| + \int \middle| gh| \right. \\
 & {\leq \parallel f\underset{p}{\parallel} \parallel h\underset{q}{\parallel} + \parallel g\underset{p}{\parallel} \parallel h\underset{q}{\parallel}} \\
 & \left. = ( \parallel f\underset{p}{\parallel} + \parallel g\underset{p}{\parallel})(\int \middle| f + g \middle| {}_{(p - 1)q})^{1/q} \right.
\end{matrix}\]

其中 \(q\) 是 \(p\) 的 Hölder conjugate. 这里的 punchline is actually: 由于

\[q: = \frac{p}{p - 1}\]

actually,

\[(p - 1)q = p\]

因而:

\[\begin{matrix}
\left. \int \middle| f + g|^{p} \right. & \left. \leq ( \parallel f\underset{p}{\parallel} + \parallel g\underset{p}{\parallel})(\int \middle| f + g \middle| {}_{(p - 1)q})^{1/q} \right. \\
 & \left. = ( \parallel f\underset{p}{\parallel} + \parallel g\underset{p}{\parallel})(\int \middle| f + g \middle| {}_{p})^{1/q} \right. \\
 & \left. = ( \parallel f\underset{p}{\parallel} + \parallel g\underset{p}{\parallel})(\int \middle| f + g \middle| {}_{p})^{1 - 1/p} \right.
\end{matrix}\]

两边同时除以 \(\left. (\int \middle| f + g \middle| {}_{p})^{1 - 1/p} \right.\) 得到:

\[\left. (\int \middle| f + g \middle| {}_{p})^{1/p} = : \parallel f + g\underset{p}{\parallel} \leq \parallel f\underset{p}{\parallel} + \parallel g\underset{p}{\parallel} \right.\]

从而得证.

\end{proof}

\begin{qlremark}{Remark}

这里的技巧是: 把一个 \(p\) 次方的函数拆成一个 \(1\) 次方的函数和一个 \(p - 1\) 次方的函数, 并且使用Hölder, 这样就得到了一个 1 次的函数的 \(p\)-norm 和另一个 \(p - 1\) 次的函数的 \(q\) norm, 但是注意 \(q(p - 1) = p\), 因而这个函数就变成了

\[\left. (\int \middle| \phi \middle| {}_{p})^{1/q} \right.\]

的形式. 并且注意到:

\[\left. \frac{\left. \int \middle| \phi|^{p} \right.}{\left. (\int \middle| \phi \middle| {}_{p})^{1/q} \right.} = (\int \middle| \phi \middle| {}_{p})^{1/p} = \parallel \phi\underset{p}{\parallel} \right.\]

\end{qlremark}

\begin{qlremark}{Remark}

Minkowski 不等式证明的是 \(1 \leq p < \infty\) 时的 \(p\)-norm 的三角不等式. 但是对于 \(0 < p < 1\), 它并不成立. 因为这个时候 \(p - 1 < 0\), 我们刚才的证明不作效.\\
直观的证明: 在 \(p \geq 1\) 的时候, \(|x|^{p}\) 是一个 strictly convex 的函数; 而在 \(0 < p < 1\) 的时候, \(|x|^{p}\) 则是一个 strictly concave 的函数.\\
因而我们运用 strictly concave 的性质:

\[\left. |a + b \middle| {}_{p} > \middle| a \middle| {}_{p} + \middle| b|^{p} \right.\]

再由积分可得到反例. (比如取 indicator function 进行积分)

\end{qlremark}

\subsection{\texorpdfstring{properties of \(L^{p}\) spaces (\(1 \leq p < \infty\))}{properties of L\^{}\{p\} spaces (1 \textbackslash leq p \textless{} \textbackslash infty)}}\label{properties-of-lp-spaces-1-leq-p-infty}

\subsection{\texorpdfstring{\(L^{p}\) (\(1 \leq p < \infty\)) is Banach}{L\^{}\{p\} (1 \textbackslash leq p \textless{} \textbackslash infty) is Banach}}\label{lp-1-leq-p-infty-is-banach}

\begin{theorem}{\protect\phantomsection\label{kn-l-p-space-1-leq-p-infty-is-banach}{\textbf{\(L^{p}\)
space (\(1 \leq p < \infty\)) is Banach}}}

\(L^{p}\) (\(1 \leq p < \infty\)) is Banach.

\end{theorem}

\begin{proof}

By last lec 的定理: 一个 NVS 是 Banach 的等价条件是任意 abs conv series 都 conv. 因而我们证明这一点即可.\\
Suppose \(f_{n} \in L^{p}\) for each \(n\), 并且这个 series abs conv, 即:

\[B := \sum\limits_{k = 1}^{\infty} \parallel f_{k}\underset{p}{\parallel} < \infty\]

我们 define:

\[g(x): = \sum\limits_{k = 1}^{\infty}f_{k}(x),\quad g_{n}(x): = \sum\limits_{k = 1}^{n}f_{k}(x)\]

我们 WTS:

\[\lim\limits_{n\rightarrow\infty}g_{n} = g\]

in \(p\)-norm induced metric sense, 即, for some \(f \in L^{p}\), 有

\[\lim\limits_{n\rightarrow\infty} \parallel g - g_{n}\underset{p}{\parallel} = 0\]

我们 Set:

\[\left. G_{n} := \sum\limits_{k = 1}^{n} \middle| f_{k} \middle| ,\quad G := \sum\limits_{k = 1}^{\infty} \middle| f_{k}| \right.\]

这个函数以及函数列的定义是为了使用 DCT, 作 donimating function 用.\\
By measurable function 的 limit behavior, 有

\[G_{n},G \in L^{+}\]

并且

\[\parallel G_{n}\underset{p}{\parallel} \leq \sum\limits_{k = 1}^{n} \parallel f_{k}\underset{p}{\parallel} \leq B\]

由于 \(G_{n}\operatorname{\nearrow ︎}G\), by MCT 有

\[\int G^{p} = \lim\limits_{n\rightarrow\infty}\int G_{n}^{p} \leq B^{p} < \infty\]

由于 \(G \in L^{p}\), 有

\[G(x) < \infty\quad a.e.\]

于是:

\[g(x): = \sum\limits_{k = 1}^{\infty}f_{k}(x) < \infty\quad a.e.\]

又 \(\left. |g_{n} \middle| , \middle| g \middle| \leq G,g_{n}\rightarrow g \right.\), 可得到:

\[\left. |g_{n} - g \middle| {}_{p} \leq 2^{p}G^{p} \in L^{1} \right.\]

因而 by DCT 可以得到:

\[\left. \lim\limits_{n}\int \middle| g_{n} - g \middle| {}_{p} = 0 \right.\]

从而

\[\left. \lim\limits_{n\rightarrow\infty} \parallel g - g_{n}\underset{p}{\parallel} = (\lim\limits_{n}\int \middle| g_{n} - g \middle| {}_{p})^{1/p} = 0 \right.\]

\end{proof}

\begin{qlremark}{Remark}

1. 我们说一个 function seq converge to 一个 function 指的是 in the sense of distance, 而这里就是 metric induced by norm, 即\textbf{它们的差的 \(L_{p}\) norm converge to \(0\).}\\
2. 注意, 我们 recall: \(f_{k}\rightarrow f\) a.e. 并不说明 \(f_{k}\rightarrow f\) in \(L^{1}\), 因为每个点 converge 的速度不一样. 当然, 对 \(L^{p}\) 也同理.\\
3. \textbf{虽然 a.e. convergence 不能推出 \(L^{p}\) convergence, 但是配合 DCT, 则可以推出.} \textbf{DCT 是我们证明 \(L^{p}\) convergence 的关键.}\\
4. 要证明

\[\lim\limits_{n\rightarrow\infty} \parallel g - g_{n}\underset{p}{\parallel} = 0\]

完全可以忽略积分外的 \(1/p\) 次方. 其实只需要证明

\[\left. \lim\limits_{n}\int \middle| g_{n} - g \middle| {}_{p} = 0 \right.\]

就可以了. 证明 \(L^{p}\) convergence, 比起 \(L^{1}\) convergence 略困难的地方就是被积函数变得更大了.

\end{qlremark}

\subsection{\texorpdfstring{Criterion for \(L^{p}\) convergence: 逐点 a.e. conv \(+\) \(L^{p}\) 积分值 conv}{Criterion for L\^{}\{p\} convergence: 逐点 a.e. conv + L\^{}\{p\} 积分值 conv}}\label{criterion-for-lp-convergence-ux9010ux70b9-a.e.-conv-lp-ux79efux5206ux503c-conv}

我们刚才 mention: DCT 对于 function seq \(L^{p}\) convergence 的证明有很大作用. 这里我们就提供一个 DCT 推出的 \(L^{p}\) convergence 的判断准则:

\begin{theorem}{\protect\phantomsection\label{kn-criterion-for-l-p-convergence}{\textbf{Criterion
for \(L^{p}\) convergence}}}

if \(f_{n}\rightarrow f\) a.e. and \(\parallel f_{n}\underset{p}{\parallel}\rightarrow \parallel f\underset{p}{\parallel}\), then \(\parallel f_{n} - f\underset{p}{\parallel}\rightarrow 0\).

\end{theorem}

即

\[\text{a.e. conv} + L^{p}\ \text{norm conv}\Longrightarrow L^{p}conv\]

但是 converse 并不成立. 反例是 typewriter function.

\begin{proof}

In Hw 8.

\end{proof}

\subsection{\texorpdfstring{dense subsets of \(L^{p}\), and specially \(L^{p}({\mathbb{R}},m)\)}{dense subsets of L\^{}\{p\}, and specially L\^{}\{p\}(\{\textbackslash mathbb\{R\}\},m)}}\label{dense-subsets-of-lp-and-specially-lpmathbbrm}

\begin{proposition}

对于任意 \(1 \leq p < \infty\), the set of \(\{\)simple functions\(\}\), is dense in \(L^{p}\).\\
即:

\[\left\{ {f:X\rightarrow{\mathbb{C}} \mid f = \sum\limits_{1}^{n}a_{j}\chi_{E_{j}},\mu(E_{j}) < \infty} \right\}\]

是 \(L^{p}\) 的 dense subset.

\end{proposition}

\begin{qlremark}{Remark}

我们已经 proved this for \(L^{1}\), 而其实这个 density 推广至 \(L^{p}\) 也成立.

\end{qlremark}

\begin{proof}

对 \(f\) 使用 simple function seq 逼近, 使用 \(\left. 2^{p} \middle| f|^{p} \right.\) 作为 dominating function of \(|f_{k} - f|^{p}\); 而后使用 DCT 得证.

\end{proof}

\begin{theorem}{\protect\phantomsection\label{kn-density-of-compactly-supported-continuous-functions-in-lp}{\textbf{\(C_{c}^{0}({\mathbb{R}}^{n})\)
is dense in \(L^{p}({\mathbb{R}},m)\) for \(1 \leq p < \infty\)}}}

\(C_{c}^{0}({\mathbb{R}}^{n})\) is dense in \(L^{p}({\mathbb{R}},m)\) for \(1 \leq p < \infty\)

\end{theorem}

\begin{proof}

exercise. Similar to the proof for \(L^{1}\), 只需要使用加入 \(p\) power 的 function 作为 dominating function 即可.

\end{proof}

\section{\texorpdfstring{\(L^{\infty}\) space, and relationship between \(L^{p}\) spaces (\(0 \leq p \leq \infty\)) {[}Fol 6.1, finished{]}}{L\^{}\{\textbackslash infty\} space, and relationship between L\^{}\{p\} spaces (0 \textbackslash leq p \textbackslash leq \textbackslash infty) {[}Fol 6.1, finished{]}}}

对应 Folland 6.1(3), finishing 6.1.\\
我们已经完成了对 \(1 \leq p < \infty\) 的 \(L^{p}\) space 的构建. 现在, 我们来构建最后一块拼图: \(L^{\infty}\) space.

\subsection{\texorpdfstring{\(L^{\infty}\) space}{L\^{}\{\textbackslash infty\} space}}\label{linfty-space}

我们考虑这个启发式的例子:

\[X := \left\{ {1,2,\cdots,n} \right\},\quad\mathcal{A} := \mathcal{P}(X),\quad\mu = \mu_{counting}\]

于是:

\[L^{p}(\mu) = \left\{ (a_{1},\cdots,a_{n}): \parallel (a_{1},\cdots,a_{n})\underset{p}{\parallel} = (\sum \middle| a_{i} \middle| {}_{p})^{1/p} < \infty \right\} = {\mathbb{C}}^{n}\]

我们发现:

\[\left. \parallel (a_{1},\cdots,a_{n})\underset{p}{\parallel}\rightarrow\max\limits_{j} \middle| a_{j} \middle| \quad\text{as}\quad p\rightarrow\infty \right.\]

因为 \(p\) 取得越大, 最大的 entry 的 contribution 占比就越突出.\\
对于这样的 \(L^{p}\) space, 我们可以定义 \(\sup\) norm, 定义为最大的 entry.\\
即便 \(X\) 是 countable 的, 这个定义也可以定义为 \(\left. \sup_{j} \middle| a_{j}| \right.\), make sense.\\
那么如果我们想要给任意的 measure space 定义 sup norm 呢? 我们可以考虑

\[\left. \parallel f\underset{\infty}{\parallel}: = \sup\limits_{x \in X} \middle| f(x) \middle| ? \right.\]

实际上我们有更好的定义方式:

\begin{definition}{\protect\phantomsection\label{kn-essential-supremum}{\textbf{essential
supremum}}}

\[\parallel f\underset{\infty}{\parallel}: = \inf\left\{ {a \geq 0:\mu\left\{ x: \middle| f(x) \middle| > a \right\} = 0} \right\}\]

也可以写作:

\[\left. \text{ess}\sup\limits_{x \in X} \middle| f(x)| \right.\]

\end{definition}

\begin{qlremark}{Remark}

essential sup 是一个比较容易搞错的定义.\\
一个 function 的 essential supremum 即: 这个 function 几乎处处的 sup.\\
它 \(\leq \sup f\) , 因为它允许在零测集上存在一些点的函数值大于它.\\
这是合理的, 因为积分可以不考虑零测集.\\

\end{qlremark}

\begin{qlremark}{Remark}

对于零测集只有空集的 measure space 上的函数, 比如 对于 \(\ell^{\infty}({\mathbb{N}})\) 上的函数, 其 essentail supermum 即 supermum.\\
对于

\[\left. \sup\limits_{x \in X} \middle| f(x)| \right.\]

我们也有一个称呼, 称其为 \textbf{uniform norm}. 即:

\[\left. \parallel f\underset{u}{\parallel} := \sup\limits_{x \in X} \middle| f(x)| \right.\]

\end{qlremark}

\begin{definition}{\protect\phantomsection\label{kn-l-infty-space}{\textbf{\(L^{\infty}\)
space}}}

\[L^{\infty}(\mu): = \left\{ {f:X\rightarrow{\mathbb{C}}\ \text{measurable}: \parallel f\underset{\infty}{\parallel} < \infty} \right\}/ \sim\]

where \(\sim\) 表示 a.e. 相等的函数的 equiv class.

\end{definition}

\begin{qlremark}{Remark}

注意: \textbf{\(f \in L^{\infty}(\mu)\), 并不等价于 \(f\) a.e. bounded!}\\
实则 recall: \(f\) a.e. \textbf{bounded 是 \(f \in L^{p}(\mu)\) for any \(1 \leq p \leq \infty\) 的必要条件}, 否则, 函数积分不可能 \(< \infty\), 函数 \(p\) 次方的积分更加不可能 \(< \infty\).\\
\(f \in L^{\infty}(\mu)\) 是一个很严格的条件, 当然严格强于 \(f\) a.e. bounded.\\
比方说: \textbf{\(f = \frac{1}{x}\), 只有在 \(0\) 这一个点上 \(f\) 是 unbounded 的, 但是它的 essential supermum 仍然是 \(\infty\)}, 因为不可能通过去掉一个 measure \(0\) set 来使它 bounded.\\
无法找到一个 \(M\), 使得 \(f\) 在几乎处处都小于 \(M\). 你只能控制, \(f\) 在 \((0,1/M)\) 上小于 \(M\), 这个集合的测度随 \(M\) 增大越来越小, 但是永远都是正测度. (同样这个函数也不属于任何 \(L^{p}(m)\).)\\
一个函数 essential supermum \(< \infty\), 即 \(\in L^{\infty}\), 则必须要它 unbounded 的这个行为是可以忽略不计的, 不能是明显的. 比如它在 \(\mathbb{Q}\) 上 unbounded. 如果是在一个点上连续 blow up, 那么它就不可能 \(\in L^{\infty}\). 类似于这里的 \(f = \frac{1}{x}\).\\

\end{qlremark}

\begin{qlremark}{Remark}

我们在本节课还会证明, 如果 measure space \(X\) has finite measure, 那么有

\[L^{\infty}(X) \subset \cdots \subset L^{p}(X) \subset \cdots \subset L^{q}(X) \subset \cdots \subset L^{1}(X)\]

for 任意的 \(p \geq q\).\\
这表明的是, 在一定要求下, \(L^{\infty}\) 是要求最严格的 space.

\end{qlremark}

下面是一个比较典型的例子:

\subsection{\texorpdfstring{\(\ell^{\infty}\) space}{\textbackslash ell\^{}\{\textbackslash infty\} space}}\label{ellinfty-space}

\begin{definition}{\protect\phantomsection\label{kn-ell-infty}{\textbf{\(\ell^{\infty}\)}}}

\[\ell^{\infty}: = \left\{ (a_{j})_{1}^{\infty}: \parallel (a_{j})\underset{\infty}{\parallel}: = \sup\limits_{j} \middle| a_{j} \middle| < \infty \right\}\]

\end{definition}

\begin{example}

\[f = x\chi_{\mathbb{Q}} \in L^{\infty}(m)\]

with

\[\parallel f\underset{\infty}{\parallel} = 0\]

因为整个 \(\mathbb{Q}\) 都是零测的.

\end{example}

\begin{qlremark}{Remark}

\(\ell^{\infty}\) 其实就是:

\[X: = {\mathbb{N}},\quad\mathcal{A}: = \mathcal{P}(X),\quad\mu = \mu_{counting}\]

的 measure space 上的 \(L^{\infty}(\mu)\).

\[\ell^{\infty} = L^{\infty}({\mathbb{N}},\mathcal{P}({\mathbb{N}}),\mu_{counting})\]

一个 seq 就是一个从 \(\mathbb{N}\) to \(\mathbb{C}\) 的函数, 把每个 entry map to 一个 complex number.\\
\textbf{而对于 counting measure 作为 measure 的 measure space 上, 唯一的零测集就是空集}, 因为哪怕只取一个元素, 这个子集的测度也是 1.\\
比如, 我们只取三个 entry \(1,2,8\), 看 \(\left\{ {|a_{n}|} \right\}_{1}^{\infty}\backslash\left\{ |a_{1} \middle| , \middle| a_{2} \middle| , \middle| a_{3} \middle| ) \right\}\) 中的 \(\sup\) value, 也不符合 essential supremum 的定义.\\
因而我们发现, \textbf{对于 唯一的零测集就是空集 的 measure space, for example, 任何以 counting measure 作为 measure 的 measure space, 其 essential sup norm 就是普通的 sup value norm.}\\
比如

\[{\mathbb{C}}^{1},{\mathbb{C}}^{2},{\mathbb{C}}^{3},\cdots,\ell^{\infty}\]

\end{qlremark}

\subsection{\texorpdfstring{\(L^{\infty}\) 的基本性质: as a NVS; Hölder's ineq on it; dense subsets}{L\^{}\{\textbackslash infty\} 的基本性质: as a NVS; Hölder's ineq on it; dense subsets}}\label{linfty-ux7684ux57faux672cux6027ux8d28-as-a-nvs-huxf6lders-ineq-on-it-dense-subsets}

\begin{lemma}

如果 \(f \in L^{\infty}(\mu)\) 则:

\begin{itemize}
\item
  一定有 \(\left. |f(x) \middle| \leq \parallel f\underset{\infty}{\parallel} \right.\) for a.e. \(x\).
\item
  存在一个 bounded 函数 \(g\), 使得 \(f = g\) a.e.
\end{itemize}

\end{lemma}

\begin{proof}

显然.

\end{proof}

\begin{qlremark}{Remark}

是否有在某个零测集上 unbounded 但是却 \(L^{\infty}\) 的函数? 答案是肯定的:

\[f(x) = \left\{ \begin{matrix}
{\frac{1}{x},} & {x \in {\mathbb{Q}} \cap (0,1\rbrack} \\
{0,} & \text{otherwise}
\end{matrix} \right.\]

有 \(\parallel f\underset{\infty}{\parallel} = 0\).

\end{qlremark}

\begin{theorem}

\begin{itemize}
\item
  \[\parallel fg\underset{\infty}{\parallel} \leq \parallel f\underset{1}{\parallel} \parallel g\underset{\infty}{\parallel}\]

  可以把它看作 \textbf{Hölder 的一部分特殊情况}, 因为可以看作

  \[\frac{1}{1} + \frac{1}{\infty} = 1\]

  从而补充完整了 Hölder ineq for \(1 \leq p,q \leq \infty\)
\item
  \(L^{\infty}\) 是一个 \textbf{normed vector space}, equipped with \(\parallel \cdot \underset{\infty}{\parallel}\)
\item
  simple functions are dense in \(L^{\infty}\)
\end{itemize}

\end{theorem}

\begin{proof}

容易证明.

\end{proof}

\begin{qlremark}{Remark}

注意, \(L^{\infty}\) 和 \(L^{p}\) 有一个出入点是: \textbf{\(C_{c}^{0}({\mathbb{R}}^{n})\) 并不是 \(L^{\infty}({\mathbb{R}}^{n},m)\) 上的 dense subspace!}\\

\end{qlremark}

\subsection{\texorpdfstring{\(L^{\infty}\)-convergence 作为 (finite measure space 下) 最强的 \(L^{p}\) convergence: 等价于 uni. conv a.e.}{L\^{}\{\textbackslash infty\}-convergence 作为 (finite measure space 下) 最强的 L\^{}\{p\} convergence: 等价于 uni. conv a.e.}}\label{linfty-convergence-ux4f5cux4e3a-finite-measure-space-ux4e0b-ux6700ux5f3aux7684-lp-convergence-ux7b49ux4ef7ux4e8e-uni.-conv-a.e.}

\begin{theorem}{\protect\phantomsection\label{kn-convergence-in-l-infty-iffuniform-convergence-a-e}{\textbf{convergence
in \(L^{\infty}\) \(\Leftrightarrow\)uniform convergence a.e.}}}

\[f_{n}\rightarrow f\text{in}\ L^{\infty}\Leftrightarrow\text{exists null set}\ E \subset X s.t.f_{n}\rightarrow f\ \text{uniformly on}\ E^{c}\]

(注意, 这\textbf{不是 conv almost uniformly}, 而是一个比 almost uniformly \textbf{更强}的条件: \textbf{conv uniformly almost everywhere}, 因为 almost uniformly 只要求对于任意的 \(\epsilon\), 都存在一个 measure 小于 \(\epsilon\) 的 \(E\), 使得在 \(E^{c}\) 上 uni conv 即可.)

\end{theorem}

\begin{qlremark}{Remark}

这一条 convergence 十分惊人. 因为\textbf{对于普通的 \(L^{p}\) space, converge in \(L^{p}\) 和 a.e. convergence 并没有任何的互推关系}; 但是对于 \(L^{\infty}\) convergence, 我们却可以把它\textbf{等价于 uniform convergence almost everywhere}, which is 一个\textbf{比 a.u convergence 更强, 比 a.e. convergence 更强的逐点 convergence}. 可以看出 \(L^{\infty}\) convergence 是比任何 \(L^{p}\) convergence 都要强一个层次的收敛性质.\\
这一点

\end{qlremark}

\begin{proof}

⇐: Suppose \(f_{n}\rightarrow f\) uni. a.e; WTS: \(f_{n}\rightarrow f\) in \(L^{\infty}\) \(f_{n}\rightarrow f\) uni. a.e 即: 存在零测集 \(E \subset X\), \(f_{n}\rightarrow f\) on \(E^{c}\).\\
Let \(\epsilon > 0\).\\
\(f_{n}\rightarrow f\) uni. a.e 表明, 存在 \(N\) 使得 for all \(n \geq N\) 有

\[\left. \forall x \in E^{c},\quad \middle| f_{n}(x) - f(x) \middle| < \epsilon \right.\]

by def, exactly is:

\[\left. \parallel f_{n} - f\underset{L^{\infty}}{\parallel} = \text{ess\,sup}_{x \in X} \middle| f_{n}(x) - f(x) \middle| \leq \epsilon \right.\]

This shows that \(\parallel f_{n} - f\underset{L^{\infty}}{\parallel}\rightarrow 0\), 即 \(f_{n}\rightarrow f\) in \(L^{\infty}\).\\
⇐: Suppose \(f_{n}\rightarrow f\) in \(L^{\infty}\); WTS: \(f_{n}\rightarrow f\) uni. a.e.Denote:

\[\epsilon_{n} := \parallel f_{n} - f\underset{L^{\infty}}{\parallel}\]

By assumption, \(\epsilon_{n}\rightarrow 0\). Define for each \(n\):

\[A_{n} := \left\{ x \in X: \middle| f_{n}(x) - f(x) \middle| > \epsilon_{n} \right\}\]

By def \(\left. \parallel f_{n} - f\underset{L^{\infty}}{\parallel} = \text{ess sup}_{x} \middle| f_{n}(x) - f(x) \middle| \leq \epsilon_{n} \right.\), 于是 \(\mu(A_{n}) = 0\) 那么令:

\[E := \bigcup\limits_{n = 1}^{\infty}A_{n}\]

by subadditivity of measure 有 \(\mu(E) = 0\). 于是对于任意 \(\epsilon_{n}\), 都有

\[\left. |f_{n}(x) - f(x) \middle| \leq \epsilon_{n}\rightarrow 0,\quad\text{for all}\ x \in E^{c} \right.\]

由于 \(\epsilon_{n}\rightarrow 0\), showing that outside \(E\), 有 \(\parallel f_{n} - f\underset{L^{\infty}}{\parallel}\rightarrow 0\).

\end{proof}

\begin{qlremark}{Remark}

这两个 convergence 直觉上是自然相等的.\\
但是这并不能够说明 \(L^{\infty}(\mu)\text{-convergence}\) 就是强于任何 \(L^{p}(\mu)\text{-convergence}\) 的. 因为即便是 uniform 的 ptwise conv 也无法推出 \(L^{p}\) conv.\\
特殊情况是, 如果整个 base space \(X\) 是 finite measure 的, 则可以推出

\[L^{\infty}(\mu)\text{-convergence}\Longrightarrow L^{p}(\mu)\text{-convergence}\Longrightarrow L^{q}(\mu)\text{-convergence}\Longrightarrow\cdots\]

whenever \(p > q\). (可证明)\\
但是对于无限测度空间, 这种推论未必成立.

\end{qlremark}

\subsection{\texorpdfstring{\(L^{\infty}\) as Banach space}{L\^{}\{\textbackslash infty\} as Banach space}}\label{linfty-as-banach-space}

\begin{theorem}{\protect\phantomsection\label{kn-l-p-1-leq-p-leq-infty-is-banach}{\textbf{\(L^{p}\)
(\(1 \leq p \leq \infty\)) is Banach}}}

For any measure space \((X,\mathcal{A},\mu)\), \(L^{p}(\mu)\) is Banach for all \(1 \leq p \leq \infty\)

\end{theorem}

\begin{proof}

我们已经 proved 了 \(1 \leq p < \infty\) 的 case, 现在 prove \(p = \infty\) 的 case.\\
By \hyperref[thm-09-l-p-space-and-inequalities-another-criterion-for-banach-space]{Theorem~\ref*{thm-09-l-p-space-and-inequalities-another-criterion-for-banach-space}}, 我们知道 STS: every abs conv series conv in \(L^{\infty}\).\\
我们 suppose \(f_{k} \in L^{\infty}\) 有

\[\sum\limits_{k = 1}^{\infty} \parallel f_{k}\underset{\infty}{\parallel} < \infty\]

WTS: \(\sum_{k = 1}^{\infty}f_{k}\) converges.\\
Set:

\[E_{k} := \left\{ x: \middle| f_{k}(x) \middle| > \parallel f_{k} \middle| |_{\infty} \right\}\]

于是有

\[\mu(E_{k}) = 0\quad\text{for each}\ k\]

因而 setting

\[E: = \bigcup\limits_{k = 1}^{\infty}E_{k}\]

有

\[\mu(E) = 0\]

note:

\[\left. x \in E^{c}\Longrightarrow\sum\limits_{k = 1}^{\infty} \middle| f_{k}(x) \middle| \leq \sum\limits_{k = 1}^{\infty} \parallel f_{k}\underset{\infty}{\parallel} < \infty \right.\]

从而,

\[g: = \sum\limits_{k = 1}^{\infty}f_{k}\]

在 \(E^{c}\) 上是 well-defined 的, 且 bounded by \(\sum_{k = 1}^{\infty} \parallel f_{k}\underset{\infty}{\parallel}\).\\
对于 \(x \in E\), 我们可以随便设置值, 比如 \(\pi\), 然后 define \(g(x) = \pi\) on \(x \in E\). 然后对于 each \(n\), 我们 set:

\[\begin{matrix}
{g_{n}(x): = \{\sum\limits_{k = 1}^{n}f_{k}(x),\quad x \in E^{c}} \\
{\frac{1}{\pi},\quad x \in E}
\end{matrix}\]

从而

\[\begin{matrix}
\left. \parallel g_{n} - g\underset{\infty}{\parallel} \leq \sup\limits_{x \in E^{c}} \middle| g_{n}(x) - g(x)| \right. & \left. \leq \sup\limits_{x \in E^{c}} \middle| \sum\limits_{n + 1}^{\infty}f_{k}(x)| \right. \\
 & \left. \leq \sup\limits_{x \in E^{c}}\sum\limits_{n + 1}^{\infty} \middle| f_{k}(x)| \right. \\
 & {\leq \sum\limits_{n + 1}^{\infty} \parallel f_{k}\underset{\infty}{\parallel}\rightarrow 0}
\end{matrix}\]

\end{proof}

\begin{qlremark}{Remark}

My reflection: 不 Banach 的 normed vector space 是什么样子的呢? 即, 这个 space 中存在某些 series, 其对应的 norm series absolutely conv 但是它却不 converge to 一个元素呢?\\
我们考虑空间 \(c_{00}\)，它是所有 finite supp 的 seq 组成的空间:

\[c_{00} := \left\{ {x = (x_{1},x_{2},\ldots) \in {\mathbb{R}}^{\mathbb{N}} \mid \text{only finite}\ x_{i} \neq 0} \right\}\]

with \(\ell^{1}\) norm:

\[\left. \parallel x \parallel = \sum\limits_{i = 1}^{\infty} \middle| x_{i}| \right.\]

\(c_{00}\) 是一个 normed vector space, 但不是 Banach space, 它的完备化是 \(\ell^{1}\). 我们考虑 series, with:

\[x_{n} = e_{n}/2^{n}\]

其中 \(e_{n} = (0,\ldots,0,1,0,\ldots)\), 第 \(n\) 个位置是 1, 其余是 \(0\), 是这个 NVS 的 standard basis.\\
显然每个 \(x_{n} \in c_{00}\)，并且：

\[\parallel x_{n} \parallel = \frac{1}{2^{n}}\quad\Rightarrow\quad\sum\limits_{n = 1}^{\infty} \parallel x_{n} \parallel = \sum\limits_{n = 1}^{\infty}\frac{1}{2^{n}} = 1\]

这是一个 absolutely convergent series, 但其和

\[\sum\limits_{n = 1}^{\infty}x_{n} = \left( {\frac{1}{2},\frac{1}{4},\frac{1}{8},\ldots} \right) \notin c_{00}\]

\(L^{p}\) space 的 Banach 性表示了其\textbf{极限存在的稳定性}. recall, Banach 即 complete NVS, 而 \textbf{complete 是比 closed 更强的条件}.\\
因而\textbf{任何一个 \(L^{p}\) 函数列, 如果 Cauchy / converge in \(L^{p}\) norm, 那么它的极限一定在 \(L^{p}\) 里.}

\end{qlremark}

\subsection{\texorpdfstring{relationship between \(L^{p}\) spaces}{relationship between L\^{}\{p\} spaces}}\label{relationship-between-lp-spaces}

\subsection{\texorpdfstring{\(L^{m}(\mu) \subset L^{n}(\mu),0 < n \leq m \leq \infty\), for measure finite space)}{L\^{}\{m\}(\textbackslash mu) \textbackslash subset L\^{}\{n\}(\textbackslash mu),0 \textless{} n \textbackslash leq m \textbackslash leq \textbackslash infty, for measure finite space)}}\label{lmmu-subset-lnmu0-n-leq-m-leq-infty-for-measure-finite-space}

刚才我们已经 state 了, 但还没有证明:

\begin{theorem}{\protect\phantomsection\label{kn-inclusion-relation-between-l-p-spaces-when-base-space-is-finite}{\textbf{inclusion
relation between \(L^{p}\) spaces (when base space is finite measure)}}}

如果 measure space \(X\) has finite measure, 那么有

\[L^{\infty}(X) \subset \cdots \subset L^{m}(X) \subset \cdots \subset L^{n}(X) \subset \cdots\]

for 任意的 \(m \geq n\).

\end{theorem}

这是我们首次把 \(p < 1\) 也 include 进我们的讨论.

这个 statement 即: 对于 from finite measure space to \(\mathbb{C}\) 的 function \(f\), 它的 \(\parallel f\underset{m}{\parallel} < \infty\) 是比 \(\parallel f\underset{n}{\parallel} < \infty\) 更强的条件.\\
尤其, 除去 \(L^{\infty}\) 的情况, 它更直接的意思是: 对于 \(0 < n \leq m < \infty\) 而言, \(f\) 的绝对值的 \(m\) 次方的积分 \(< \infty\) 是比 \(f\) 的绝对值的 \(n\) 次方的积分 \(< \infty\) 要更强的条件.

这其实是一件比较直观的事情. 因为对于 \(\left. |f \middle| \geq 1 \right.\) 的部分,

\[\left. \int_{|f| \geq 1} \middle| f \middle| {}_{large} \geq \int_{|f| \geq 1} \middle| f|^{small} \right.\]

而对于 \(\left. |f \middle| < 1 \right.\) 的部分,

\[\left. \int_{|f| < 1} \middle| f \middle| {}_{large} \leq \int_{|f| < 1} \middle| f|^{small} \right.\]

然而\textbf{由于整个 space 的 measure 是 finite 的, \(\left. |f \middle| < 1 \right.\) 的部分并不影响}. 因为

\[\left. \int_{|f| < 1} \middle| f \middle| {}_{large} \leq \int_{|f| < 1} \middle| f \middle| {}_{small} \leq \int_{|f| < 1}1 \leq \mu(X) \right.\]

因而, 对于 \(\mu(X) < \infty\) 的情况, 显然有 \(\parallel f\underset{large}{\parallel} < \infty\) 是比 \(\parallel f\underset{small}{\parallel} < \infty\) 更强的条件.\\
\textbf{(实际上, 如果只有 measure finite 的 \(x\) 上 \(\left. |f(x) \middle| < 1 \right.\), 那么即便 \(\mu(X) = \infty\), \(\parallel f\underset{large}{\parallel} < \infty\) 也是比 \(\parallel f\underset{small}{\parallel} < \infty\) 更强的条件; 而如果有 measure infinite 的 \(x\) 上 \(\left. |f(x) \middle| < 1 \right.\), 那么有可能 \(\parallel f\underset{large}{\parallel} < \infty\) 是比 \(\parallel f\underset{small}{\parallel} < \infty\) 更弱的条件)}\\
My point: 虽然说 \(|f(x)|^{large}\) 比起 \(|f(x)|^{small}\) 是更大还是更小取决于 \(|f(x)|\) 是否 \(\geq 1\) or \(< 1\), 但是 \(\geq 1\) 的值是可以 unbounded 的, 而 \(< 1\) 的值再怎么通过小次方变得更大, 也超不过 \(1\). 因而 \(\left. |f(x) \middle| \geq 1 \right.\) 的部分通常更能函数积分值的有限性, 除非在一个 measure infinite 的集合上 \(\left. |f(x) \middle| < 1 \right.\).

这里有一个更加严格的证明:

\begin{proof}

首先, 对于 \(m = \infty\) 的 case, 如果 \(f \in L^{m} = L^{\infty}\), 那么取任意 \(1 \leq n < \infty\) 都有:

\[\left. \int \middle| f \middle| {}_{n} \leq \int \parallel f\underset{\infty}{\overset{n}{\parallel}} = \parallel f\underset{\infty}{\overset{n}{\parallel}}\mu(X) < \infty \right.\]

其次, 对于正常的 \(m < \infty\) 的 case, 我们使用 Hölder: 如果 \(f \in L^{m}\), 那么对于任意 \(n < m\), 我们可以构造出 Hölder conjugate \(\frac{m}{n}\) 和 \(\frac{m}{m - n}\),从而:

\[\begin{matrix}
\left. \int \middle| f|^{n} \right. & \left. = \int \middle| f \middle| {}_{n} \cdot 1 \right. \\
 & \left. \leq (\int( \middle| f \middle| {}_{n})^{\frac{m}{n}})^{\frac{n}{m}}(\int 1^{\frac{m}{m - n}})^{\frac{m - n}{m}} \right. \\
 & {= \parallel f\underset{m}{\overset{n}{\parallel}}\mu(X)^{\frac{m - n}{m}} < \infty}
\end{matrix}\]

从而

\[\parallel f\underset{m}{\parallel} < \infty\Longrightarrow \parallel f\underset{n}{\parallel} < \infty\]

这一 proof 利用 Hölder conjuate, 通过构造包含 \(\frac{m}{n}\) 的 Hölder conjugate, 把 \(\left. \int \middle| f|^{n} \right.\) 改成了 \(\parallel f\underset{m}{\parallel}\) 的 expression.

\end{proof}

以下是一个经典的例子:

\begin{example}

考虑 \textbf{measure finite 的 measure space \((0,1)\)}: 通过经典的 Calculus 我们知道:

\[f(x) = \frac{1}{x^{m}} \in L^{p}(0,1)\quad\text{for all}\ p < \frac{1}{m}\]

但是对于任意的 \(m\), 都有:

\[f(x) = \frac{1}{x^{m}} \notin L^{p}(0,1)\]

而我们再看一个 \textbf{measure infinite 的 measure space \((1,\infty)\) 上的反例}, 采用同一个函数:

\[f(x) = \frac{1}{x^{m}},\quad x \in (1,\infty)\]

这个时候, \(p\) 越大, \(\left. \int \middle| f \middle| {}_{p} = \parallel f\underset{p}{\overset{p}{\parallel}} \right.\) 反而越小, 通过经典的 Calculus 我们知道:我们知道 而对于

\[f(x) = \frac{1}{x^{m}} \in L^{p}(1,\infty)\quad\text{for all}\ p > \frac{1}{m}\]

并且 \(f \in L^{\infty}(1,\infty)\), 因为 \(\parallel f\underset{\infty}{\parallel} = 1\).\\
这个空间上的这个函数正对应了我们刚才讨论的, 如果有 infinite measure 数量的 \(x\) 上 \(\left. |f(x) \middle| < 1 \right.\), 那么很可能 \(\parallel f\underset{large}{\parallel} < \infty\) 是比 \(\parallel f\underset{small}{\parallel} < \infty\) 更弱的条件

\end{example}

\subsection{\texorpdfstring{control arbitrary \(\parallel f\underset{m}{\parallel}\) 和 \(\parallel f\underset{n}{\parallel}\) 的大小比例, in measure finite space}{control arbitrary \textbackslash parallel f\textbackslash underset\{m\}\{\textbackslash parallel\} 和 \textbackslash parallel f\textbackslash underset\{n\}\{\textbackslash parallel\} 的大小比例, in measure finite space}}\label{control-arbitrary-parallel-fundersetmparallel-ux548c-parallel-fundersetnparallel-ux7684ux5927ux5c0fux6bd4ux4f8b-in-measure-finite-space}

\begin{qlremark}{Remark}

刚才我们的推导中,

\[\left. \int \middle| f \middle| {}_{n} \leq \parallel f\underset{m}{\overset{n}{\parallel}}\mu(X)^{\frac{m - n}{m}} < \infty \right.\]

两边开 \(p\) 方, 可以得到一个不等式:

\begin{theorem}

对于 measure finite space \(X\), 对于任意的 \(0 < n \leq m \leq \infty\), 有:

\[\parallel f\underset{n}{\parallel} \leq \parallel f\underset{m}{\parallel}\,\mu(X)^{\frac{1}{n} - \frac{1}{m}}\]

\end{theorem}

这也是一个有用的不等式. 它在 measure finite space 上, 对于任意的可测函数, 控制了两个任意的 function \(p\)-norm (虽然 for \(p < 1\) 不能严格地称为 norm) 之间的大小关系。

\end{qlremark}

\subsection{\texorpdfstring{\((L^{n} \cap L^{r}) \subset L^{m} \subset (L^{n} + L^{r})\), 对任意 \(0 < n < m < r \leq \infty\)}{(L\^{}\{n\} \textbackslash cap L\^{}\{r\}) \textbackslash subset L\^{}\{m\} \textbackslash subset (L\^{}\{n\} + L\^{}\{r\}), 对任意 0 \textless{} n \textless{} m \textless{} r \textbackslash leq \textbackslash infty}}\label{ln-cap-lr-subset-lm-subset-ln-lr-ux5bf9ux4efbux610f-0-n-m-r-leq-infty}

\begin{proposition}

对于 measurable \(f:X\rightarrow{\mathbb{C}}\),

\[t\mapsto \parallel f\underset{\frac{1}{t}}{\parallel}\]

is \textbf{log-convex}.\\
equivalently 即: 对于任意的 \(0 < n < m < r \leq \infty\), 都有

\[\parallel f\underset{m}{\parallel} \leq \parallel f\underset{n}{\overset{\lambda}{\parallel}} \cdot \parallel f\underset{r}{\overset{1 - \lambda}{\parallel}}\]

where

\[\lambda := \frac{\frac{1}{m} - \frac{1}{r}}{\frac{1}{n} - \frac{1}{r}} \in (0,1),\quad i.e.(\frac{1}{m}) = \lambda(\frac{1}{n}) + (1 - \lambda)(\frac{1}{r})\]

\end{proposition}

\begin{qlremark}{Remark}

log convex 即: 这个函数的 \(\log\) 函数是 convex 的. 即对于任意 \(x,y\), 以及 \(\lbrack x,y\rbrack\) 上的任意一点, 即 \(\lambda x + (1 - \lambda)y\) for some \(\lambda \in \lbrack 0,1\rbrack\), 都有:

\[\log f(\lambda x + (1 - \lambda)y) \leq \lambda\log f(x) + (1 - \lambda)\log f(y)\]

即:

\[f(\lambda x + (1 - \lambda)y) \leq f(x)^{\lambda}f(y)^{1 - \lambda}\]

例如: \(e^{x},e^{x^{2}},x^{x}\) 都是 log-convex 的. convex 函数的几何意义是 \textbf{"函数值小于等于两端的线性插值"}, 中点值 \(\leq\)两端值的\textbf{算术平均}, 而 log-convex 函数的几何意义是: , 中点值 \(\leq\)两端值的\textbf{几何平均}.\\
这里, 两端点是 \(\frac{1}{r} < \frac{1}{n}\), 而中间的取点则是 \(\frac{1}{m}\). log convexity 性质表明:

\[\parallel f\underset{m}{\parallel} \leq \parallel f\underset{n}{\overset{\lambda}{\parallel}} \cdot \parallel f\underset{r}{\overset{1 - \lambda}{\parallel}}\]

\end{qlremark}

\begin{proof}

For \(r = \infty\), then \(\lambda = \frac{n}{m}\).\\
Since

\[\left. |f \middle| {}_{m} = \middle| f \middle| {}_{n} \cdot \middle| f \middle| {}_{m - n} \leq \middle| f \middle| {}_{n} \cdot \parallel f\underset{\infty}{\overset{m - n}{\parallel}}\quad a.e. \right.\]

可以得到

\[\left. \int \middle| f \middle| {}_{m} \leq (\int \middle| f \middle| {}_{n}) \cdot \parallel f\underset{\infty}{\overset{m - n}{\parallel}} = \parallel f \middle| |_{n}^{n} \cdot \parallel f\underset{\infty}{\overset{m - n}{\parallel}} \right.\]

从而 Taking \(q\)th root 得到结果:

\[\parallel f\underset{m}{\parallel} \leq \parallel f\underset{n}{\overset{n/m}{\parallel}} \parallel f\underset{\infty}{\overset{1 - n/m}{\parallel}}\]

For \(r < \infty\): 我们采用 conjugate exponents:

\[\frac{n}{\lambda m},\frac{r}{(1 - \lambda)m}\]

这是因为:

\[(\frac{1}{m}) = \lambda(\frac{1}{n}) + (1 - \lambda)(\frac{1}{r})\Longrightarrow 1 = \lambda(\frac{m}{n}) + (1 - \lambda)(\frac{m}{r})\]

从而 Applying Hölder:

\[\begin{matrix}
\left. \int \middle| f|^{m} \right. & \left. = \int \middle| f \middle| {}_{\lambda m} \middle| f|^{(1 - \lambda)m} \right. \\
 & \left. \leq (\int \middle| f \middle| {}_{n})^{\frac{\lambda m}{n}}(\int \middle| f \middle| {}_{r})^{\frac{(1 - \lambda)m}{r}} \right. \\
 & {= \parallel f\underset{n}{\overset{\lambda m}{\parallel}} \cdot \parallel f\underset{r}{\overset{(1 - r)m}{\parallel}}}
\end{matrix}\]

Taking \(q\) th root 得到结果.

\end{proof}

\begin{qlremark}{Remark}

Hölder's ineq 仍然是这里重要的一步. 我们这里需要利用 convexity 表述中的 "point on a line segment" 条件来构造一个 conjugate.

\end{qlremark}

\begin{qlremark}{Remark}

此处我们可以由这个 proposition 直接得到一个推论: ~

\begin{corollary}

对于任意的 \(0 < n < m < r \leq \infty\), 都有

\[(L^{n} \cap L^{r}) \subset L^{m}\]

\end{corollary}

\end{qlremark}

\begin{example}

令 \(A\) 为任意集合, \(0 \leq p < q \leq \infty\), 有:

\[\parallel f\underset{q}{\parallel} \leq \parallel f\underset{p}{\parallel}\quad\text{and thus}\quad\ell^{p}(A) \subset \ell^{q}(A)\]

这是因为

\[\left. \parallel f\underset{\infty}{\overset{p}{\parallel}} = \sup\limits_{\alpha} \middle| f(\alpha) \middle| {}_{p} \leq \sum\limits_{\alpha} \middle| f(\alpha) \middle| {}_{p} = \parallel f\underset{p}{\overset{p}{\parallel}} \right.\]

于是 for \(q \neq \infty\) case

\[\parallel f\underset{q}{\parallel} \leq \parallel f\underset{p}{\overset{\lambda}{\parallel}} \parallel f\underset{\infty}{\overset{1 - \lambda}{\parallel}} \leq \parallel f\underset{p}{\parallel}\]

(另一 case, trivial.)\\
我们发现 \(\ell^{p}\) 空间, \(p\) 越小要求反而越严格.\\
这是因为 \(\ell^{p}\) 空间中一个函数就是一个 seq, 其 \(p\)-norm 就是各项的 \(p\) 次方和, 再开 \(p\) 次方根.\\
对于一个 seq, 如果它的累和 series 收敛, 它的各项肯定是 \textbf{eventually 收敛的}, 那么这些除了有限项外的这些项的绝对值都是 \(< 1\) 的, 那么 \textbf{\(p\) 越大, 它们 \(p\) 次方和只会越小}. 这正对应了我们之前说的 \textbf{"\(\left. |f(x) \middle| < 1 \right.\) 的点主导函数" 的情况.}\\
\strut \\

\end{example}

相对于这个inclusion 关系, 我们还有另外一个 inclusion 关系:

\begin{proposition}{\protect\phantomsection\label{kn-l-m-l-n-l-r-0-n-m-r-leq-infty}{\textbf{每个
\(L^{m}\) 函数都是一个 \(L^{n}\) 函数和一个 \(L^{r}\) 函数的和
(\(0 < n < m < r \leq \infty\))}}}

对于任意的 \(0 < n < m < r \leq \infty\), 都有

\[L^{m} \subset (L^{n} + L^{r})\]

\end{proposition}

这个 inclusion 关系有一种调和的感觉在里面. 它 roughly mean 给定一个函数, 它可以拆成一个更加容易积的函数和一个更加不容易积的函数, 并且我们很大程度上可以控制这两个函数的可积性.\\
但其实很简单, 就是用我们之前的 \(\left. |f(x) \middle| < 1 \right.\) 和 \(\geq 1\) 的点作为区分, 把函数的定义域分成两部分. 如果 \(|f|\) 的 m 次方是可积的, 那么更小的 \(n\) 次方, 对于 \(\left. |f(x) \middle| \geq 1 \right.\) 的部分肯定也是可积的; 更大的 \(r\) 次方, 对于 \(\left. |f(x) \middle| < 1 \right.\) 的部分肯定也是可积的;

\begin{proof}

Suppose \(f \in L^{m}\). Let

\[E := \left\{ x: \middle| f(x) \middle| > 1 \right\}\]

let

\[g: = f\chi_{E},\quad h := f\chi_{E^{c}}\]

于是 \(g \in L^{n}\) for all \(0 < n \leq m\), \(h \in L^{r}\) for all \(r \geq m\) and \(r = \infty\).

\end{proof}
\end{document}
